% Generated draft: requires mathematical and source review.
\documentclass[11pt]{article}
\usepackage[margin=25mm]{geometry}
\usepackage{fontspec}
\setmainfont{Latin Modern Roman}
\usepackage{amsmath,amssymb,mathtools,unicode-math}
\setmathfont{Latin Modern Math}
\usepackage{enumitem,needspace,xurl,hyperref}
\setlength{\parindent}{0pt}
\setlength{\parskip}{5pt}
\setlength{\emergencystretch}{3em}
\newcommand{\N}{\mathbb{N}}
\newcommand{\Z}{\mathbb{Z}}
\newcommand{\Q}{\mathbb{Q}}
\newcommand{\R}{\mathbb{R}}
\newcommand{\C}{\mathbb{C}}
\newcommand{\F}{\mathbb{F}}
\newcommand{\A}{\mathbb{A}}
\newcommand{\PP}{\mathbb{P}}
\newcommand{\E}{\mathbb{E}}
\newcommand{\eps}{\varepsilon}
\newcommand{\dd}{\,\mathrm{d}}
\begin{document}
\section*{UnsolvedMath ID: 20001050}
\textbf{Source problem code:} catalogue entry
\subsection{Definitions and mathematical statement}
The literal record is only ``Conjecture 6.9. Under the hypotheses of''. Neither the antecedent nor conclusion is present. The following vertexwise-load statement is recovered by the attached research aid and remains a qualified reconstruction despite the registry's \texttt{exact} label.

Let $D$ be a nonempty finite simple loopless digraph on $n$ vertices, and let each ordered arc $vw$ have a finite label set $S_{vw}$ (empty sets can be treated as absent arcs). Write
\[
 \ell(v)=\sum_{w\ne v}|S_{vw}|.
\]
A rainbow directed cycle admits distinct representatives from its arc label sets, including a digon when present. The recovered Conjecture 6.9 asserts that $\ell(v)\geq n$ for every $v$ forces a rainbow cycle.

\textbf{Reported cut-critical reduction.} If a counterexample exists, the aid proves that deleting vertices and label incidences yields a counterexample on $m\geq2$ vertices with exact load
\[
 \ell(v)=m\qquad\text{for every vertex},
\]
and, for every nonempty proper $X\subset V$, the strict cut certificate
\[
 \max_{v\notin X}\sum_{x\in X}|S_{vx}|\geq|X|+1.
\]
The maximum is existential in the outside vertex; the assertion is not that every outside vertex has this load into $X$. All operations retain a rainbow-cycle-free instance.

In particular the overlap core, consisting of arcs with $|S_{vw}|\geq2$, is strongly connected. On every directed cycle in that core, Hall's distinct-representative condition fails: there is a subset $F$ of at least three cycle arcs with
\[
 \left|\bigcup_{e\in F}S_e\right|<|F|.
\]
The lower size three follows because each core arc already has two available labels. This gives necessary structure for hypothetical counterexamples, not an example proving they exist.

\textbf{Reported threshold and hypothesis distinctions.} A directed $m$-cycle whose arcs all share $m-1$ labels has load $m-1$ everywhere and no rainbow cycle. Thus the pointwise threshold $m$ cannot simply be replaced by $m-1$.

For layers $E_i=\{vw:i\in S_{vw}\}$, the aid distinguishes the stronger hypothesis of the preceding recovered Conjecture 6.8,
\[
 \sum_i\min_v d^+_{E_i}(v)\geq n,
\]
from the more permissive vertexwise condition
\[
 \min_v\sum_i d^+_{E_i}(v)\geq n.
\]
The first implies the second, but the report supplies a directed-cycle family showing no converse; the exact family is not printed in the summary. Proving a result only under the sum-of-minima hypothesis therefore does not prove the full vertexwise conjecture. The cut-critical theorem leaves the latter unresolved.
\subsection{Short English statement}
The source conjecture is truncated. For the aid's vertexwise rainbow-cycle reconstruction, any counterexample has a smaller exact-load kernel with strict incoming-cut certificates, a strongly connected multi-label core, and Hall obstructions on its cycles. These are conditional restrictions, not a proof or a counterexample to the conjecture.
\subsection{Sources}
\begin{itemize}
\item Local UnsolvedMath registry, permanent ID $20001050$, \texttt{catalogue entry}; canonical \texttt{aim-combinatorics-notes.json}, \texttt{notes[174]}.
\item AIM workshop ``The Caccetta--Haggkvist conjecture,'' PDF fragment Conjecture 6.9: \url{https://aimath.org/WWN/caccetta/caccetta.pdf}. The verification calls the assigned text an extraction fragment; the neighboring Conjecture 6.8 is discussed only through the attached summary.
\item Research attempt 1, \texttt{partial\_result}, run status \texttt{open}, internal rigor \texttt{proved}, low novelty confidence; current registry partially solved. Explicitly machine-generated, not peer reviewed, and to be independently verified before citation. No complete external references or original literature entry supplied.
\end{itemize}
\end{document}
